\documentclass[11pt]{letter}
\usepackage[margin=1in]{geometry}
\usepackage{hyperref}

\signature{Ian Todd\\Sydney Medical School\\University of Sydney\\itod2305@uni.sydney.edu.au}
\address{Ian Todd\\Sydney Medical School\\University of Sydney\\Sydney, NSW 2006\\Australia}

\begin{document}

\begin{letter}{The Editors\\Physical Review X\\American Physical Society}

\opening{Dear Editors,}

I am pleased to submit our manuscript, ``The Dimensional Landauer Bound: Why Information Flow Requires Geometric Work,'' for consideration as a Research Article in \textit{Physical Review X}.

\textbf{Summary.} Landauer's principle gave physics a profound insight: erasing one bit of information has a minimum energetic cost of $k_B T \ln 2$. This result anchors the thermodynamics of computation in logical irreversibility. However, modern physical and biological information-processing systems---brains, learning algorithms, and complex dynamical substrates---do not primarily operate on discrete bits. They operate on high-dimensional manifolds: continuous neural population states, field configurations, latent representations in machine learning, and quantum states in Hilbert space.

In this work, we derive the ``missing half'' of Landauer's principle. While Landauer's bound describes the one-time cost of \textit{erasing} a bit, we identify the continuous metabolic cost of \textit{maintaining} a low-dimensional representation against the eroding force of noise---effectively the ``rent'' biology must pay for structure. Whenever a physical system compresses high-dimensional dynamics into a lower-dimensional representational channel, there is an unavoidable energetic cost that we call \textit{dimensional work}. This cost arises from geometric irreversibility, not logical irreversibility.

\textbf{Main Result.} We derive a Dimensional Landauer Bound of the form
\[
W_{\min} \ge k_B T \ln 2\, \Delta I + k_B T\, C_\Phi,
\]
where $\Delta I$ is the classical Landauer term (bits erased) and $C_\Phi$ is a dimensional contraction cost defined via KL divergence and connected geometrically to the Jacobian of the projection. This leads to a physically transparent statement:
\begin{itemize}
\item Landauer's principle quantifies the cost of erasing bits.
\item The Dimensional Landauer Bound quantifies the cost of representing information on a lower-dimensional manifold.
\end{itemize}

\textbf{Methods and Numerical Illustration.} Using stochastic thermodynamics, optimal control, and information geometry, we derive the bound in a general setting of high-dimensional Langevin dynamics subject to dimensional projection. We instantiate the bound numerically in four complementary settings:
\begin{enumerate}
\item 2D Brownian motion on curved 1D manifolds (geometric work scales with curvature);
\item Kuramoto oscillators (coherence reduces control work for equivalent readout fidelity);
\item Autoencoder training under SGD-as-Langevin (effort explosion below intrinsic dimensionality);
\item Neural sheet with ephaptic coupling (energy-matched random vs.\ coherent fields isolates the geometric mechanism).
\end{enumerate}

\textbf{Why Physical Review X?} This manuscript sits at the intersection of statistical physics, information theory, and complex systems, with direct implications for neuroscience and machine learning---a scope aligned with PRX's mission to publish cross-cutting, high-impact physics. The Dimensional Landauer Bound provides a foundational extension of well-established thermodynamic principles into a domain---manifold-based computation and high-dimensional inference---that is increasingly central to modern physics and AI.

To our knowledge, no prior work provides a unified, thermodynamically grounded treatment of the energetic cost of dimensional reduction across physical, biological, and artificial systems. We therefore believe this manuscript will be of broad interest to the PRX audience.

We confirm that this manuscript has not been published and is not under consideration elsewhere. We have no conflicts of interest to declare.

\closing{Sincerely,}

\end{letter}
\end{document}
